\documentclass{ximera}

\title{Limits Activity Facilitator Guide}
\author{MATH 425: Calculus I}

\begin{document}
\begin{abstract}
    Working with the peers in your group, solve the following problems. Make sure to show and justify all your work. Make sure everyone in the group understands the solution and participates. Be prepared to report your answers to the whole class. 
\end{abstract}
\maketitle

%\textbf{Student Learning Outcomes}
%Students will be able to:
% \begin{itemize}   
%    \item Evaluate one- and two-sided limits using a graph.
    %\item Identify when a limit does not exist. 
%    \item Sketch a graph based on limit and function value information.
%     \item Identify where the limit of a piecewise function fails to exist and explain their answer.
%    \item Use limit laws to evaluate limits given values of limits for unspecified functions.
%    \item (Challenge) Investigate where the division limit law may fail.

%\textbf{Prerequisite Knowledge:} This activity expects that students will be familiar with limits on a graph, and the limit laws.  This activity is a mixture of the Limit Activity and the Limit Laws Activity.  \\
%\\

%This activity seems to either go quickly or expose groups' misunderstanding of limits. Encourage students to dig deeper into what a limit is/is not, as this will be the foundation of all the work we do this semester (and into Calc II). The piecewise function tends to trip students up, as they just sub in and call it a day (if they get that far).  The table problem also causes problems for students.  This is the first of many table problems in this collection, so letting students know they should expect more of this type of thinking will be helpful. 


\begin{exercise}
     If the given graph is the graph of the function $y=g(x)$, state the value of the following.  If the limit does not exist, explain why.  
\\
  \begin{image}
    \includegraphics{limit1.png}
  \end{image}
  \xmalt{A graph of a function with the following characteristics: a smooth curve leading to an open circle point at (-3,4), then a smooth curve down to a closed point at (0,1), jumping to a open point at (0,-1) and a smooth curve through an open point at (2,2) into infinite oscillation until x=5, suddenly becoming almost parabola-like at (5,3) with a maximum occuring at x=7.}
  \begin{enumerate}
    \item $\lim_{x\rightarrow -4}g(x)$ \textcolor{blue}{$=3$}
    \item $g(-4)$ \textcolor{blue}{ $=3$}
    \item $\lim_{x\rightarrow -3^-}g(x)$ \textcolor{blue}{ $=4$}
    \item $\lim_{x\rightarrow -3^+}g(x)$ \textcolor{blue}{ $=4$}
    \item $\lim_{x\rightarrow -3}g(x)$ \textcolor{blue}{ $=4$}
    \item $g(-3)$ \textcolor{blue}{ undefined}
    \item $\lim_{x\rightarrow 0}g(x)$ \textcolor{blue}{ does not exist}
    \item $g(0)$ \textcolor{blue}{ $=1$}
    \item $\lim_{x\rightarrow 2^+}g(x)$ \textcolor{blue}{ $=2$}
    \item $\lim_{x\rightarrow 2^-}g(x)$ \textcolor{blue}{ $=2$}
    \item $\lim_{x\rightarrow 2}g(x)$ \textcolor{blue}{ $=2$}
    \item $\lim_{x\rightarrow 5^-}g(x)$ \textcolor{blue}{ undefined}
    \item $\lim_{x\rightarrow 5^+}g(x)$ \textcolor{blue}{ $=3$}
    \item $\lim_{x\rightarrow -2}g(x)$ \textcolor{blue}{ $=2$}
    \item $g(-2)$ \textcolor{blue}{ $=2$}
   
  \end{enumerate}
\end{exercise}

\begin{exercise}
    On a coordinate plane, draw a function, $f$, that satisfies the following conditions.  Be neat in your drawing, and label any important points.
        \begin{enumerate}
            \item $\displaystyle\lim_{x\rightarrow -1^+}f(x)=-2$
            \item $\displaystyle\lim_{x\rightarrow -1^-} f(x)=3$
            \item $f(x)$ has domain $[-5, -1)\cup(-1,7] $
            \item $\displaystyle\lim_{x\rightarrow 2}f(x)=0$
            %\item $f(x)$ has a removable discontinuity at $x=2$.
            %\item $f(x)$ is left continuous at $x=7$.
            %\item $\displaystyle\lim_{x\rightarrow 4}f(x)=-\infty$.
            \item $\displaystyle\lim_{x\rightarrow -3}f(x)=4$
            \item $\displaystyle\lim_{x\rightarrow 7^-} f(x)=-3=f(7)$
            \item $f(2)\neq 0$
        \end{enumerate}
        \textcolor{blue}{ Here is a possible solution:
        \begin{image}
            \includegraphics{limits2A.jpg}
        \end{image}}
        \xmalt{A graph with the following properties: a smooth curve connecting a closed point at (-5,-2) to an open point at (-1,3), a jump to an open point at (-1,-2), followed by a smooth curve with an open point at (2,0) with a singular point at (2,2), and a continued curve to a closed point at (7,-3)}
        %\textcolor{orange}{This type of problem is a struggle for some students.  Make sure students know that we want all of these conditions to be true in the graph of a single function, not a graph for each condition.  If students need more support, suggest that they draw the ``local'' behavior and try to connect it after.  This is also a great time to remind students to check whether their graph actually is a function.}

\end{exercise}



\begin{exercise}
    Where does the limit of $f(x)$ fail to exist?  
      \[f(x)=
    \begin{cases}
    3x-1 &\text{if }x<-2\\
    2x^2 &\text{if} -2\leq x\leq 1\\
    2x+7 & \text{if } 1<x
    \end{cases}\]
    \textcolor{blue}{ $\displaystyle \lim_{x\rightarrow -2^-}f(x)=\lim_{x\rightarrow -2^-} 3x-1=3(-2)-1=-7$\\
    $\displaystyle \lim_{x\rightarrow -2^+} f(x)=\lim_{x\rightarrow -2^+ }2x^2=2(-2)^2=8$\\
    Thus, $\displaystyle \lim_{x\rightarrow -2}f(x)$ DNE\\
    $\displaystyle\lim_{x\rightarrow 1^-} f(x)=\lim_{x\rightarrow 1^-} 2x^2=2(1)^2=2$\\
    $\displaystyle\lim_{x\rightarrow 1^+} f(x)=\lim_{x\rightarrow 1^+} 2x+7=2(1)+7=9$\\
    Therefore $\displaystyle\lim_{x\rightarrow 1} f(x)$ DNE\\
    On each interval, the rules are polynomials, so the limit exists.  The only places the limit fails to exist are at $x=-2$ and $x=1$} \\
    %\textcolor{orange}{Make sure students are not just subbing numbers into the two functions and comparing them.  We want to highlight how the limits are being used here.}
\end{exercise}

\begin{exercise}
    Use the given table to evaluate the following limits, or explain why you cannot. \\
  \begin{center}
        \begin{tabular}{c|c|c}
            $c$ & $\lim_{x\rightarrow c}f(x)$ & $\lim_{x\rightarrow c}g(x)$  \\ \hline
            -2 &  5 & 3\\ 
            0 & -1 & DNE\\ 
            2 & $\displaystyle\frac{1}{2}$ & -2\\
            3 & 0 & 4
        \end{tabular}
     \end{center}
     \begin{enumerate}
     \item $\displaystyle\lim_{x\rightarrow -2} f(x)(g(x))^2$ \textcolor{blue}{ $\displaystyle=\lim_{x\rightarrow -2}f(x)(\lim_{x\rightarrow -2} g(x))^2=5(3)^2=45$}
     \item $\displaystyle\lim_{x\rightarrow 0} \frac{g(x)}{f(x)}$ \textcolor{blue}{ $\displaystyle=\frac{\lim_{x\rightarrow 0}g(x)}{\lim_{x\rightarrow 0}f(x)}$ cannot be calculated, since $\lim_{x\rightarrow 0}g(x)$ does not exist}
     \item $\displaystyle \lim_{x\rightarrow 0} \sqrt[3]{f(x)}$ \textcolor{blue}{ $\displaystyle=\sqrt[3]{\lim_{x\rightarrow 0}f(x)}=\sqrt[3]{-1}=-1$}
     \item $\displaystyle \lim_{x\rightarrow 2} 3f(x)-2g(x)+7x$ \textcolor{blue}{ $\displaystyle=3\lim_{x\rightarrow2}f(x)-2\lim_{x\rightarrow 2}g(x)+7\lim_{x\rightarrow2}x=2(\frac{1}{2})-2(-2)+7(2)=1+4+14=19$}
     \item $\displaystyle \lim_{x\rightarrow 3} \frac{2f(x)+1}{\sqrt{g(x)}}$ \textcolor{blue}{ $\displaystyle=\frac{2\lim_{x\rightarrow 3}f(x)+\lim_{x\rightarrow3} 1}{\sqrt{\lim_{x\rightarrow3} g(x)}}=\frac{2(0)+1}{\sqrt{4}}=\frac{1}{2}$}

    \end{enumerate}
\end{exercise}

\begin{exercise}
     \textbf{Challenge problem: }Find an example where $\displaystyle \lim_{x\rightarrow 0} \frac{f(x)}{g(x)}$ exists, but neither $\displaystyle\lim_{x\rightarrow 0}f(x)$ or $\displaystyle \lim_{x\rightarrow 0}g(x)$ exist.\\
     \textcolor{blue}{ Suppose $f(x)=\frac{1}{x^3}$ and $g(x)=\frac{1}{x^5}$.  Then $\lim_{x\rightarrow 0}f(x)$ and $\lim_{x\rightarrow 0}g(x)$ DNE.\\
 $\displaystyle\lim_{x\rightarrow 0}\frac{f(x)}{g(x)}=\lim_{x\rightarrow 0}\frac{\frac{1}{x^3}}{\frac{1}{x^5}}=\lim_{x\rightarrow 0}\frac{x^5}{x^3}=\lim_{x\rightarrow 0}x^2=0$ }\\
 %\textcolor{orange}{This is a difficult question that really asks students to reason as to why we need the hypotheses of limits existing for the limit laws to apply. We do not necessarily call out this fact, so if a group is really getting into this one, it may be worth throwing out the question, ``Does this counter the quotient law for limits?  Why not?''}
\end{exercise}



\end{document}